% =============================================================================
%  LE: Level Emergence — skeleton paper
%  STATUS: UNPROVEN. Independent of S1/N1/CAT1; can be pursued in parallel.
% =============================================================================
\documentclass[11pt]{article}

\usepackage[utf8]{inputenc}
\usepackage[T1]{fontenc}
\usepackage[expansion=false]{microtype}
\usepackage{amsmath, amssymb, amsthm, mathtools}
\usepackage[margin=1in]{geometry}
\usepackage{enumitem}
\usepackage[hidelinks]{hyperref}

\theoremstyle{plain}
\newtheorem{theorem}{Theorem}
\newtheorem{lemma}[theorem]{Lemma}
\newtheorem*{namedconjecture}{Conjecture}
\theoremstyle{definition}
\newtheorem{definition}[theorem]{Definition}

\title{Level Emergence: The Four-Condition Mechanism by which \\
       Hierarchical Life Arises Dynamically \\
       \large (UNPROVEN --- skeleton paper for LE)}
\author{Liam McGuire\thanks{Unaffiliated researcher, Seattle, USA. liammcg44@gmail.com}}
\date{\today}

\begin{document}
\maketitle

\begin{abstract}
This paper is a \textbf{skeleton} that records LE from the
scaffolding paper~\cite{emergent_systems_2026}. \textbf{No proof is
attempted.} LE is the framework's structural theory of major
evolutionary transitions, generalised from Maynard Smith \&
Szathm\'ary~\cite{maynardsmith1995major} to a substrate-agnostic
setting. It is the dynamical companion to the static
characterisation theorems (S1, N1, CAT1). The result is UNPROVEN.
\end{abstract}

\section{Statement (UNPROVEN)}

\begin{quote}
\emph{``LE (level emergence): a system $\mathcal{S}$ at vitality
profile depth $k$ undergoes a level transition to depth $k+1$ over a
time window iff its level-$k$ entities satisfy four conditions over
that window: persistence, aggregation, group variation, and group
selection.''}
\end{quote}
\noindent (Scaffolding paper~\cite{emergent_systems_2026},
\S\texttt{sec:le}; restated in \S\texttt{sec:le-conjecture}.)

\section{Plan of attack (no proof attempted)}

Four sub-lemmas, one per condition, plus the iff bridging them.

\begin{enumerate}[leftmargin=2em,itemsep=4pt]
  \item \textbf{Persistence sub-lemma (UNPROVEN).} Level-$k$ entities
    have lifetimes much longer than substrate relaxation time. Make
    precise what ``much longer'' means in the operadic setting.
  \item \textbf{Aggregation sub-lemma (UNPROVEN).} A non-trivial
    coupling structure $\kappa^{(k)}$ produces persistent multi-entity
    patterns. Connect to the substrate composition operator
    $\boxtimes_\kappa$.
  \item \textbf{Group-variation sub-lemma (UNPROVEN).} The coupling
    varies non-trivially across the population of aggregates.
    Operationalise as positive Shannon entropy of $\kappa^{(k)}$
    across aggregate types.
  \item \textbf{Group-selection sub-lemma (UNPROVEN).} Aggregate
    types have measurably different persistence rates.
    Operationalise as positive coefficient of variation across
    persistence rates.
  \item \textbf{Necessity direction (UNPROVEN, easier).} If a level
    transition has occurred, the four conditions held over a
    preceding window. Empirically grounded in
    Maynard Smith \& Szathm\'ary~\cite{maynardsmith1995major} and
    Buss~\cite{buss1987evolution}; formalising for substrate-agnostic
    systems is the work.
  \item \textbf{Sufficiency direction (UNPROVEN, harder).} If the
    four conditions hold over a sufficient time window, a level
    transition occurs. Requires either a substrate-agnostic
    construction or a case analysis across operadic generators.
\end{enumerate}

\section{Worked specialisations}

The scaffolding paper walks through chemistry $\to$ cells, cells $\to$
organisms, organisms $\to$ ecosystems. Each is a candidate
specialisation of LE. Other candidates worth working out:

\begin{itemize}[leftmargin=1.5em,itemsep=2pt]
  \item Tierra-arms-race level-1 $\to$ level-2 transition (programmatic);
  \item spoken $\to$ written language transition (cultural);
  \item endosymbiosis-as-level-transition (Margulis, biological).
\end{itemize}

\section{Status and follow-up}

\textbf{Current status: UNPROVEN.} Independent of S1/N1/CAT1; can be
pursued in parallel. Coordinate with `c2_hierarchical_viability/`
(C2 is the post-transition structure).

\bibliographystyle{plain}
\bibliography{references}

\end{document}
